% fluxograma da metodologia. cada caixa aponta o script ou comando do repositorio
\begin{tikzpicture}[
  node distance=4mm,
  caixa/.style={draw, rounded corners=2pt, align=center, font=\scriptsize, minimum width=7.6cm, inner sep=3pt},
  dado/.style={draw, align=center, font=\scriptsize, minimum width=7.6cm, inner sep=3pt, fill=black!6},
  seta/.style={-{Latex[length=2mm]}, semithick}]
\node[dado] (base) {FashionVerse-MME: 126.928 registros em 31.333 looks (JSON)};
\node[caixa, below=of base] (pre) {Pré-processamento: junta registros repetidos (102.370 peças),\\cor Lab, divisão dos looks em treino, validação e teste};
\node[caixa, below=of pre] (atr) {Atributos do texto (estilo, ocasião, formalidade, estampa)\\e 12 grupos de cor com k-means};
\node[caixa, below=of atr] (est) {Estruturas: KD-tree, grafo peça--look (BFS de 2 passos),\\Union-Find e Kruskal (grupos), tabelas de pontos};
\node[caixa, below=of est] (ver) {Versões A, B, C, D e F: nota para todas as peças\\da categoria da peça escondida};
\node[caixa, below=of ver] (val) {Validação: escolhe pesos e parâmetros\\(metade A escolhe, metade B confirma)};
\node[caixa, below=of val] (tes) {Teste único: Hit@1, Hit@5, Hit@10 e FITB,\\10 ordens de desempate e bootstrap};
\node[dado, below=of tes] (sai) {Tabelas, gráficos de tempo, Gephi e articulações (Tarjan)};
\foreach \a/\b in {base/pre, pre/atr, atr/est, est/ver, ver/val, val/tes, tes/sai} \draw[seta] (\a) -- (\b);
\end{tikzpicture}
